% Partie : set-theory
% Chapitre : ordinals
% Section : opps

\documentclass[../../../include/open-logic-section]{subfiles}

\begin{document}

\olfileid[fr]{sth}{ordinals}{opps}
\olsection{Ordinaux successeurs et ordinaux limites}

Nous utiliserons beaucoup les ordinaux dans les prochains chapitres.
Il sera donc utile d'introduire quelques notions simples.

\begin{defn}
Pour tout ordinal $\alpha$, son \emph{successeur} est
$\ordsucc{\alpha} =\alpha \cup \{\alpha\}$. Nous disons que
$\alpha$ est un ordinal \emph{successeur} s'il existe un ordinal
$\beta$ tel que $\ordsucc{\beta} = \alpha$. Nous disons que
$\alpha$ est un ordinal \emph{limite} si et seulement s'il n'est ni
vide ni successeur.
\end{defn}
\noindent
Le résultat suivant montre que cette notion de \emph{successeur} est
la bonne :

\begin{prop}
Pour tout ordinal $\alpha$ :
\begin{enumerate}
	\item $\alpha \in \ordsucc{\alpha}$ ;
	\item $\ordsucc{\alpha}$ est un ordinal ;
	\item il n'existe aucun ordinal $\beta$ tel que
	$\alpha \in \beta \in \ordsucc{\alpha}$.
\end{enumerate}
\end{prop}

\begin{proof}
Il est immédiat que
$\alpha \in \alpha \cup \{\alpha\} = \ordsucc{\alpha}$.
De même, $\ordsucc{\alpha}$ est un ensemble transitif d'ordinaux,
donc un ordinal par le \olref[basic]{corordtransitiveord}.
Enfin, on ne peut avoir
$\alpha \in \beta \in \ordsucc{\alpha}$ : dans ce cas, soit
$\beta \in \alpha$, soit $\beta = \alpha$, en contradiction avec
le \olref[basic]{ordordered}.
\end{proof}
\noindent
Cela permet également une variante de la démonstration par récurrence
transfinie :

\begin{thm}[Récurrence transfinie simple]\ollabel{simpletransrecursion}
Soit $\phi(x)$ une formule telle que :
\begin{enumerate}
	\item $\phi(\emptyset)$ ;
	\item pour tout ordinal $\alpha$, si $\phi(\alpha)$, alors
	$\phi(\ordsucc{\alpha})$ ;
	\item si $\alpha$ est un ordinal limite et si
	$(\forall \beta \in \alpha)\phi(\beta)$, alors $\phi(\alpha)$.
\end{enumerate}
Alors tout ordinal $\alpha$ vérifie $\phi(\alpha)$.
\end{thm}

\begin{proof}
Supposons qu'un ordinal $\delta$ ne vérifie pas $\phi$. Par séparation,
les éléments de $\ordsucc{\delta}$ qui ne vérifient pas $\phi$ forment
un ensemble non vide, puisqu'il contient $\delta$. Soit $\gamma$ son
plus petit élément. Tout ordinal $\beta \in \gamma$ vérifie donc
$\phi(\beta)$.

Si $\gamma = \emptyset$, la première hypothèse donne $\phi(\gamma)$.
Si $\gamma = \ordsucc{\alpha}$, alors $\alpha \in \gamma$, donc
$\phi(\alpha)$ est vraie et la deuxième hypothèse donne $\phi(\gamma)$.
Si $\gamma$ est un ordinal limite, la troisième hypothèse donne
$\phi(\gamma)$. Dans chacun des trois cas, on obtient une contradiction.
\end{proof}
\noindent
Une dernière notation nous sera utile plus loin :

\begin{defn}\ollabel{defsupstrict}
Si $X$ est un ensemble d'ordinaux, alors
$\supstrict(X) = \bigcup_{\alpha \in X} \ordsucc{\alpha}$.
\end{defn}
\noindent
Ici, «~lsub~» signifie «~plus petite borne supérieure stricte~».
\footnote{Certains ouvrages écrivent «~$\text{sup}(X)$~» pour cette
notion. D'autres réservent toutefois «~$\text{sup}(X)$~» à la plus
petite borne supérieure \emph{non stricte}, c'est-à-dire simplement
$\bigcup X$. Si $X$ possède un plus grand élément $\alpha$, les
deux notions divergent : la plus petite borne supérieure
\emph{stricte} est $\ordsucc{\alpha}$, tandis que la plus petite
borne supérieure \emph{non stricte} est $\alpha$.} Le résultat
suivant explique pourquoi :

\begin{prop}
Si $X$ est un ensemble d'ordinaux, $\supstrict(X)$ est le plus petit
ordinal strictement supérieur à tous les ordinaux de $X$.
\end{prop}

\begin{proof}
Posons $U = (\bigcup X) \cup X$. Si $\alpha \in X$, alors
$\alpha \subseteq \bigcup X$ et $\{\alpha\} \subseteq X$, donc
$\ordsucc{\alpha} \subseteq U$. Par séparation dans l'ensemble
des parties de $U$, nous obtenons l'ensemble
$Y = \Setabs{\ordsucc{\alpha}}{\alpha \in X}$. On a alors
$\supstrict(X) = \bigcup Y$. Puisque les ordinaux sont transitifs et
que tout élément d'un ordinal est un ordinal, $\supstrict(X)$ est un
ensemble transitif d'ordinaux, donc un ordinal par
le \olref[basic]{corordtransitiveord}.

Si $\alpha \in X$, alors $\ordsucc{\alpha} \in Y$ ; par conséquent,
$\ordsucc{\alpha} \subseteq \bigcup Y = \supstrict(X)$ et donc
$\alpha \in \supstrict(X)$. Ainsi, $\supstrict(X)$ est strictement
supérieur à chaque ordinal de $X$.

Réciproquement, si $\alpha \in \supstrict(X)$, alors, pour un certain
$\beta \in X$, nous avons $\alpha \in \ordsucc{\beta} \in Y$, d'où
$\alpha \leq \beta \in X$. Ainsi, $\supstrict(X)$ est la
\emph{plus petite} borne supérieure stricte de $X$.
\end{proof}

\end{document}
